\par\noindent\SourceLeaf{47}%
\hspace*{\parindent}Initio patet ex solo termino initiali in summam soluto eventurum terminum in
$\cos 0\varphi = 1$ ductum, h.\,e. novae seriei terminum initialem; ceteri enim
termini in utraque summae ex illis provenientis parte cosinum multipli alicuius
paris ipsius $\varphi$ continent. Est igitur novae seriei terminus initialis
\[
\frac{1}{2^{2m-2}} . \frac{1}{2}\,{}^{2m-1}\mathfrak{B}^{(m)}
= \frac{1}{2^{2m-1}}\,{}^{2m-1}\mathfrak{B}^{(m)}.
\]

\par\noindent Haec vero expressio, cum sit hic coefficiens binomialis inter ceteros suae
potestatis mediorum posterior, adeoque cum altero mediorum, cui aequalis est,
addendo iunctus potestatis uno gradu altioris coefficientem eiusdem indicis
producat, sitque hac ratione
\[
{}^{2m-1}\mathfrak{B}^{(m)} .  2 = {}^{2m}\mathfrak{B}^{(m)},
\]
sic poterit exhiberi: esse terminum initialem
\[
= \frac{1}{2^{2m-1}}\frac{{}^{2m}\mathfrak{B}^{(m)}}{2}.
\]

\par\noindent Ad legem terminorum reliquorum generalem enodandam sumatur\hfill\mbox{}\linebreak $\cos\varphi^{2m-1} . \cos\varphi$
terminus generalis $r^{\text{tus}}$
\[
\frac{1}{2^{2m-2}}\,{}^{2m-1}\mathfrak{B}^{(m+r)}\cos(2r+1)\varphi . \cos\varphi.
\]
Hic, cum sit
\[
\cos(2r+1)\varphi . \cos\varphi = \tfrac{1}{2}\cos 2r\varphi + \tfrac{1}{2}\cos(2r+2)\varphi,
\]
solvetur in partes:
\[
\frac{1}{2^{2m-2}} . \frac{1}{2} . {}^{2m-1}\mathfrak{B}^{(m+r)}\cos 2r\varphi
+ \frac{1}{2^{2m-2}} . \frac{1}{2} . {}^{2m-1}\mathfrak{B}^{(m+r)}\cos(2r+2)\varphi.
\]

Ex quo adparet, ex quolibet multiplo impari hac transformatione bina paria,
alterum proxime maius, alterum proxime minus dato impari enasci. Sic nonnisi
cosinus multiplorum parium in nova serie occurrunt, quorum quilibet coefficientem
ex binis seriei adsumtae coefficientibus contiguis conflatum obtinebit. Patet
enim ex transformatione in termino $r^{\text{to}}$ ipsius
\[
\cos\varphi^{2m-1} . \cos\varphi
\]
instituta, non solum huius partem primam, quae iam ante inventa fuit
\[
\frac{1}{2^{2m-1}}\,{}^{2m-1}\mathfrak{B}^{(m+r)}\cos 2r\varphi,
\]
